\begin{table}[!h]
\centering
\caption{Table B3. Core Threshold Sensitivity Analysis}
\centering
\resizebox{\ifdim\width>\linewidth\linewidth\else\width\fi}{!}{
\fontsize{10}{12}\selectfont
\begin{threeparttable}
\begin{tabular}[t]{lrrrrrrrrrrr}
\toprule
Threshold & N Areas & Inst. Share (%) & Core N & Periph N & Core Kurt. & Periph Kurt. & Core L-kurt. & Periph L-kurt. & Core Var. & Periph Var. & Var. Ratio\\
\midrule
Top-2 (15, 20) & 2 & 59.8 & 61 & 69 & 6.3876 & 62.6077 & 0.2119 & 0.7214 & 1.5463 & 2.1806 & 0.7091\\
Baseline Top-3 (10, 15, 20) & 3 & 74.2 & 77 & 53 & 25.4703 & 12.2970 & 0.3265 & 0.7176 & 2.9408 & 0.0711 & 41.3512\\
Top-4 (4, 10, 15, 20) & 4 & 78.7 & 87 & 43 & 27.7887 & 13.7552 & 0.3256 & 0.7622 & 2.6898 & 0.0664 & 40.4809\\
\bottomrule
\end{tabular}
\begin{tablenotes}[para]
\item \textit{Note: } 
\item Restricted distribution. Baseline uses top-3 areas by instrument count (10, 15, 20).
\end{tablenotes}
\end{threeparttable}}
\end{table}
